\documentclass[10pt,a4paper]{article}
\usepackage[landscape,margin=1.5cm]{geometry}
\usepackage[T1]{fontenc}
\usepackage{amsmath,amssymb}
\usepackage[normalem]{ulem}
\usepackage{xcolor,longtable,booktabs,array}
\setlength{\tabcolsep}{4pt}
\renewcommand{\arraystretch}{1.15}
\pagestyle{empty}
\setlength{\parindent}{0pt}
\setlength{\parskip}{3pt}
\begin{document}
\section*{Summary of Changes}
Each row lists one changed sentence, or several adjacent changed sentences, from the version you examined to the Editorial Version. Deleted words appear in \textcolor{red}{\sout{red strikethrough}} and inserted words in \textcolor{blue}{\uwave{blue underline}}.
\par \textit{Score}: word-level Levenshtein distance on normalized word forms. Inserting, deleting, or replacing a word costs 1, a common function word costs 0.25, and words with the same normalized form cost 0. The score of a row is the sum over its words.
\par \textit{Page}: the page number printed on the page of the Editorial Version where the change appears.
\par Rows with a score below 1 are omitted (58 rows listed, 2 omitted). British-to-American spelling changes are not listed.
\par\medskip
\begin{longtable}{>{\raggedright\arraybackslash}p{0.7cm} >{\raggedright\arraybackslash}p{10.8cm} >{\raggedright\arraybackslash}p{10.8cm} p{0.9cm} p{1.2cm}}
\toprule No. & Original & Revised & Score & Page \\\midrule \endfirsthead
\toprule No. & Original & Revised & Score & Page \\\midrule \endhead
\bottomrule \endfoot
\bottomrule \endlastfoot
\midrule \multicolumn{5}{l}{\textbf{\large Front matter}}\\
\midrule
\multicolumn{5}{l}{\textbf{(chapter opening)}}\\[2pt]
1 & Dataset distillation, in turn, synthesizes a small, informative dataset that yields downstream performance comparable to its larger counterpart. & Dataset distillation, in turn, synthesizes a small, informative dataset that \textcolor{blue}{\uwave{aims}} \textcolor{blue}{\uwave{to}} yield downstream performance comparable to its larger counterpart. & 1.25 & xv \\
\addlinespace[3pt]
\midrule \multicolumn{5}{l}{\textbf{\large 1 Introduction}}\\
\midrule
\multicolumn{5}{l}{\textbf{1.1 Large-Scale Models and Knowledge Distillation}}\\[2pt]
2 & DD in particular optimizes a synthetic set so that models trained on it \textcolor{red}{\sout{match}} full-dataset performance, capturing essential knowledge within a minimal data budget. & DD in particular optimizes a synthetic set so that models trained on it \textcolor{blue}{\uwave{approach}} full-dataset performance, capturing essential knowledge within a minimal data budget. & 1.00 & 2 \\
\addlinespace[3pt]
\multicolumn{5}{l}{\textbf{1.4 Distributional Fidelity of the Training Signal}}\\[2pt]
3 & \textit{(new text)} & \textcolor{blue}{\uwave{In}} \textcolor{blue}{\uwave{this}} \textcolor{blue}{\uwave{thesis}} \textcolor{blue}{\uwave{it}} \textcolor{blue}{\uwave{takes}} \textcolor{blue}{\uwave{two}} \textcolor{blue}{\uwave{related}} \textcolor{blue}{\uwave{forms:}} \textcolor{blue}{\uwave{for}} \textcolor{blue}{\uwave{Chapters}} \textcolor{blue}{\uwave{\#}} \textcolor{blue}{\uwave{and}} \textcolor{blue}{\uwave{\#,}} \textcolor{blue}{\uwave{coverage}} \textcolor{blue}{\uwave{of}} \textcolor{blue}{\uwave{the}} \textcolor{blue}{\uwave{useful}} \textcolor{blue}{\uwave{parts}} \textcolor{blue}{\uwave{of}} \textcolor{blue}{\uwave{the}} \textcolor{blue}{\uwave{teacher's}} \textcolor{blue}{\uwave{input}} \textcolor{blue}{\uwave{or}} \textcolor{blue}{\uwave{representation}} \textcolor{blue}{\uwave{space;}} \textcolor{blue}{\uwave{for}} \textcolor{blue}{\uwave{Chapter}} \textcolor{blue}{\uwave{\#,}} \textcolor{blue}{\uwave{matching}} \textcolor{blue}{\uwave{an}} \textcolor{blue}{\uwave{importance-weighted}} \textcolor{blue}{\uwave{source}} \textcolor{blue}{\uwave{distribution}} \textcolor{blue}{\uwave{in}} \textcolor{blue}{\uwave{an}} \textcolor{blue}{\uwave{embedding}} \textcolor{blue}{\uwave{space.}} & 26.50 & 7 \\
\addlinespace[3pt]
4 & Yet this does not guarantee coverage: the central challenge is ensuring that synthetic samples collectively capture diverse patterns and semantics of the teacher's knowledge domain, rather than merely satisfying its hard labels. & Yet this does not guarantee coverage: the central challenge is ensuring that synthetic samples collectively capture diverse patterns, and \textcolor{blue}{\uwave{ideally}} semantics, of the teacher's knowledge domain, rather than merely satisfying its hard labels. & 1.00 & 8 \\
\addlinespace[3pt]
5 & Fidelity deficit \textcolor{red}{\sout{is}} the root cause of these failures, not a symptom. \textcolor{red}{\sout{No}} algorithmic refinement \textcolor{red}{\sout{can}} recover knowledge that is absent from the training signal. & \textcolor{blue}{\uwave{We}} \textcolor{blue}{\uwave{take}} fidelity deficit \textcolor{blue}{\uwave{to}} \textcolor{blue}{\uwave{be}} the root cause of these failures, not a symptom. Algorithmic refinement \textcolor{blue}{\uwave{alone}} \textcolor{blue}{\uwave{cannot}} recover knowledge that is absent from the training signal. & 4.75 & 9 \\
\addlinespace[3pt]
\multicolumn{5}{l}{\textbf{1.5 Thesis Contributions}}\\[2pt]
6 & Across both axes, distributional fidelity \textcolor{red}{\sout{is}} the binding constraint. The thesis addresses all three settings under a \textcolor{red}{\sout{single}} design \textcolor{red}{\sout{philosophy:}} treat distributional fidelity as a first-class objective, derive practical constraints from the deployment setting, and enforce representational coverage of the distilled signal explicitly. & Across both axes, \textcolor{blue}{\uwave{this}} \textcolor{blue}{\uwave{thesis}} \textcolor{blue}{\uwave{treats}} distributional fidelity \textcolor{blue}{\uwave{as}} the binding constraint. The thesis addresses all three settings under a \textcolor{blue}{\uwave{shared}} design \textcolor{blue}{\uwave{principle:}} treat distributional fidelity as a first-class objective, derive practical constraints from the deployment setting, and enforce representational coverage of the distilled signal explicitly. & 4.50 & 10 \\
\addlinespace[3pt]
7 & In Chapter \#, DivBFKD addresses the setting in which only a small number of real training images are available and the teacher is accessible only as a black box, returning \textcolor{red}{\sout{final}} \textcolor{red}{\sout{predictions}} rather than internal representations. & In Chapter \#, DivBFKD addresses the setting in which only a small number of \textcolor{blue}{\uwave{labeled}} real training images are available and the teacher is accessible only as a black box, returning \textcolor{blue}{\uwave{class}} \textcolor{blue}{\uwave{probabilities}} rather than internal representations. & 3.00 & 10 \\
\addlinespace[3pt]
8 & This threshold acts as a quality gate, admitting only \textcolor{red}{\sout{sufficiently}} \textcolor{red}{\sout{informative}} \textcolor{red}{\sout{and}} \textcolor{red}{\sout{class-coherent}} synthetic images \textcolor{red}{\sout{into}} the \textcolor{red}{\sout{distillation}} \textcolor{red}{\sout{pipeline,}} and \textcolor{red}{\sout{directly}} counters the \textcolor{red}{\sout{semantic}} \textcolor{red}{\sout{narrowness}} that arises when few real examples are available. We evaluate across seven image classification benchmarks and reach state-of-the-art performance, with a gain of \textcolor{red}{\sout{over}} 17\% \textcolor{red}{\sout{relative}} to standard KD on CIFAR-10 and closing to within 2\% of a fully supervised student upper bound on \textcolor{red}{\sout{the}} \textcolor{red}{\sout{simpler}} \textcolor{red}{\sout{datasets.}} & This threshold acts as a quality gate, admitting only synthetic images \textcolor{blue}{\uwave{that}} the \textcolor{blue}{\uwave{teacher}} \textcolor{blue}{\uwave{classifies}} \textcolor{blue}{\uwave{confidently,}} and counters the \textcolor{blue}{\uwave{narrow}} \textcolor{blue}{\uwave{coverage}} that arises when few real examples are available. We evaluate across seven image classification benchmarks and reach state-of-the-art performance \textcolor{blue}{\uwave{among}} \textcolor{blue}{\uwave{the}} \textcolor{blue}{\uwave{compared}} \textcolor{blue}{\uwave{few-shot}} \textcolor{blue}{\uwave{methods,}} with a gain of 17\% \textcolor{blue}{\uwave{compared}} to standard KD on CIFAR-10 and closing to within \textcolor{blue}{\uwave{about}} 2\% of a fully supervised student upper bound on \textcolor{blue}{\uwave{MNIST,}} \textcolor{blue}{\uwave{FMNIST,}} \textcolor{blue}{\uwave{and}} \textcolor{blue}{\uwave{SVHN.}} & 18.00 & 10 \\
\addlinespace[3pt]
9 & In the first phase, a diverse pool of synthetic images is built using hierarchical noise injection, nonlinear \textcolor{red}{\sout{photometric}} \textcolor{red}{\sout{transforms,}} and semantic cutmixing, producing candidate images that cover a broad region of the input space. In the second phase, a \textcolor{red}{\sout{lightweight}} primer student is trained on these candidates under \textcolor{red}{\sout{a}} contrastive objective that \textcolor{red}{\sout{maximizes}} \textcolor{red}{\sout{inter-class}} \textcolor{red}{\sout{separation,}} sharpening the diversity of the pool before full distillation begins. In the third phase, the trained primer student supplies surrogate \textcolor{red}{\sout{logits}} \textcolor{red}{\sout{as}} soft \textcolor{red}{\sout{distillation}} \textcolor{red}{\sout{labels}} to supplement the teacher's hard labels, \textcolor{red}{\sout{recovering}} \textcolor{red}{\sout{much}} \textcolor{red}{\sout{of}} the information lost to black-box access. & In the first phase, a diverse pool of synthetic images is built using hierarchical noise injection, nonlinear \textcolor{blue}{\uwave{filters,}} and semantic cutmixing \textcolor{blue}{\uwave{whose}} \textcolor{blue}{\uwave{masks}} \textcolor{blue}{\uwave{are}} \textcolor{blue}{\uwave{irregular}} \textcolor{blue}{\uwave{regions}} \textcolor{blue}{\uwave{rather}} \textcolor{blue}{\uwave{than}} \textcolor{blue}{\uwave{object}} \textcolor{blue}{\uwave{segments,}} producing candidate images that cover a broad region of the input space. In the second phase, a primer student is trained on these candidates under \textcolor{blue}{\uwave{an}} \textcolor{blue}{\uwave{instance-level}} contrastive objective that \textcolor{blue}{\uwave{pushes}} \textcolor{blue}{\uwave{each}} \textcolor{blue}{\uwave{synthetic}} \textcolor{blue}{\uwave{image}} \textcolor{blue}{\uwave{away}} \textcolor{blue}{\uwave{from}} \textcolor{blue}{\uwave{the}} \textcolor{blue}{\uwave{others,}} sharpening the diversity of the pool before full distillation begins. In the third phase, the trained primer student supplies surrogate soft \textcolor{blue}{\uwave{supervision}} to supplement the teacher's hard labels, \textcolor{blue}{\uwave{partly}} \textcolor{blue}{\uwave{compensating}} \textcolor{blue}{\uwave{for}} the information lost to black-box access. & 23.00 & 10 \\
\addlinespace[3pt]
10 & Finally, in Chapter \#, TAKE addresses dataset distillation in the natural language domain, where the goal is to distill a \textcolor{red}{\sout{compact}} \textcolor{red}{\sout{subset}} \textcolor{red}{\sout{of}} \textcolor{red}{\sout{a}} large text corpus that \textcolor{red}{\sout{can}} \textcolor{red}{\sout{surrogates}} the performance of training on the full dataset. Existing text dataset distillation methods fall short on three fronts: their weighting schemes are uniform and uninformative about task-aligned importance, their objectives optimize batch by batch rather than enforcing global distributional coverage, \textcolor{red}{\sout{and}} their outputs are embeddings rather than human-readable text. & Finally, in Chapter \#, TAKE addresses dataset distillation in the natural language domain, where the goal is to distill a large text corpus \textcolor{blue}{\uwave{into}} \textcolor{blue}{\uwave{a}} \textcolor{blue}{\uwave{compact}} \textcolor{blue}{\uwave{synthetic}} \textcolor{blue}{\uwave{surrogate}} \textcolor{blue}{\uwave{corpus,}} \textcolor{blue}{\uwave{drawn}} \textcolor{blue}{\uwave{from}} \textcolor{blue}{\uwave{an}} \textcolor{blue}{\uwave{LLM-generated}} \textcolor{blue}{\uwave{pool,}} that \textcolor{blue}{\uwave{approximates}} the performance of training on the full dataset. Existing text dataset distillation methods \textcolor{blue}{\uwave{each}} fall short on \textcolor{blue}{\uwave{at}} \textcolor{blue}{\uwave{least}} \textcolor{blue}{\uwave{one}} \textcolor{blue}{\uwave{of}} three fronts: their weighting schemes are uniform and uninformative about task-aligned importance, their objectives optimize batch by batch \textcolor{blue}{\uwave{or}} \textcolor{blue}{\uwave{per}} \textcolor{blue}{\uwave{sample}} rather than enforcing global distributional coverage, \textcolor{blue}{\uwave{or}} their outputs are embeddings rather than human-readable text. & 16.25 & 10 \\
\addlinespace[3pt]
11 & DivBFKD (image, black-box few-shot) filters synthetic candidates through generation and adaptive confidence thresholds, admitting only \textcolor{red}{\sout{semantically}} \textcolor{red}{\sout{coherent,}} \textcolor{red}{\sout{informative}} examples into distillation. & DivBFKD (image, black-box few-shot) filters synthetic candidates through generation and adaptive confidence thresholds, admitting only examples \textcolor{blue}{\uwave{the}} \textcolor{blue}{\uwave{teacher}} \textcolor{blue}{\uwave{classifies}} \textcolor{blue}{\uwave{confidently}} into distillation. & 4.25 & 11 \\
\addlinespace[3pt]
12 & Together, these contributions show that KD under restricted access is feasible and can match or exceed methods that operate under \textcolor{red}{\sout{far}} less constrained conditions. & Together, these contributions show that KD under restricted access is feasible and can match or exceed \textcolor{blue}{\uwave{several}} \textcolor{blue}{\uwave{compared}} methods that operate under less constrained conditions. & 3.00 & 12 \\
\addlinespace[3pt]
\midrule \multicolumn{5}{l}{\textbf{\large 2 Background}}\\
\midrule
\multicolumn{5}{l}{\textbf{(chapter opening)}}\\[2pt]
13 & Both tracks are connected by §\#, which establishes distributional fidelity as the binding constraint: whether the compressed artifact is a student model (Model Track) or a surrogate corpus (Dataset Track), it must \textcolor{red}{\sout{faithfully}} \textcolor{red}{\sout{recover}} the original distribution. & Both tracks are connected by §\#, which establishes distributional fidelity as the binding constraint: whether the compressed artifact is a student model (Model Track) or a surrogate corpus (Dataset Track), it must \textcolor{blue}{\uwave{adequately}} \textcolor{blue}{\uwave{represent}} the original distribution. & 2.00 & 12 \\
\addlinespace[3pt]
\multicolumn{5}{l}{\textbf{2.4 Model Distillation Under Restricted Access $\rightarrow$ 2.4.1 White-Box Model Distillation}}\\[2pt]
14 & The design space collapses to response-based distillation with \textcolor{red}{\sout{hard-label}} queries \textcolor{red}{\sout{only,}} and the information bottleneck becomes the binding constraint on what can be recovered. & The design space collapses to response-based distillation with \textcolor{blue}{\uwave{output-only}} queries, and the information bottleneck becomes the binding constraint on what can be recovered. & 2.00 & 35 \\
\addlinespace[3pt]
\multicolumn{5}{l}{\textbf{2.4 Model Distillation Under Restricted Access $\rightarrow$ 2.4.2 Black-Box Model Distillation: Few-Shot}}\\[2pt]
15 & The student's only recourse is to treat the teacher as a query interface over whatever training data is available, receiving only \textcolor{red}{\sout{hard}} \textcolor{red}{\sout{labels}} \textcolor{red}{\sout{(single}} \textcolor{red}{\sout{predicted}} class \textcolor{red}{\sout{per}} \textcolor{red}{\sout{image)}} rather than \textcolor{red}{\sout{soft}} \textcolor{red}{\sout{targets}} \textcolor{red}{\sout{(probability}} \textcolor{red}{\sout{distributions}} \textcolor{red}{\sout{over}} \textcolor{red}{\sout{all}} \textcolor{red}{\sout{classes).}} & The student's only recourse is to treat the teacher as a query interface over whatever training data is available, receiving only \textcolor{blue}{\uwave{its}} \textcolor{blue}{\uwave{output}} class \textcolor{blue}{\uwave{probabilities}} rather than \textcolor{blue}{\uwave{logits}} \textcolor{blue}{\uwave{or}} \textcolor{blue}{\uwave{internal}} \textcolor{blue}{\uwave{features.}} & 11.50 & 35 \\
\addlinespace[3pt]
16 & Combined with black-box access \textcolor{red}{\sout{that}} \textcolor{red}{\sout{provides}} only \textcolor{red}{\sout{hard}} \textcolor{red}{\sout{labels,}} this creates a double bind: signal quality and data coverage are both constrained, and neither can compensate for the other. The student receives \textcolor{red}{\sout{hard}} \textcolor{red}{\sout{labels}} only \textcolor{red}{\sout{(low}} signal quality) on a narrow set of examples (low coverage), with no obvious mechanism to escape this bind. & Combined with black-box access, \textcolor{blue}{\uwave{which}} \textcolor{blue}{\uwave{exposes}} only \textcolor{blue}{\uwave{the}} \textcolor{blue}{\uwave{teacher's}} \textcolor{blue}{\uwave{output}} \textcolor{blue}{\uwave{probabilities}} \textcolor{blue}{\uwave{and}} \textcolor{blue}{\uwave{no}} \textcolor{blue}{\uwave{logits}} \textcolor{blue}{\uwave{or}} \textcolor{blue}{\uwave{internal}} \textcolor{blue}{\uwave{features,}} this creates a double bind: signal quality and data coverage are both constrained, and neither can compensate for the other. The student receives only \textcolor{blue}{\uwave{output}} \textcolor{blue}{\uwave{probabilities}} \textcolor{blue}{\uwave{(limited}} signal quality) on a narrow set of examples (low coverage), with no obvious mechanism to escape this bind. & 13.00 & 36 \\
\addlinespace[3pt]
\multicolumn{5}{l}{\textbf{2.4 Model Distillation Under Restricted Access $\rightarrow$ 2.4.3 Black-Box Model Distillation: Data-Free}}\\[2pt]
17 & This is the specific failure mode that motivates the soft-label \textcolor{red}{\sout{recovery}} component developed in Chapter \#: \textcolor{red}{\sout{restoring}} even an approximation of the soft probability vector partially repairs the distributional fidelity of the supervision signal, independent of whether the synthetic images themselves fully cover the real data manifold. & This is the specific failure mode that motivates the \textcolor{blue}{\uwave{surrogate}} soft-label component developed in Chapter \#: \textcolor{blue}{\uwave{supplying}} even an approximation of the soft probability vector partially repairs the distributional fidelity of the supervision signal, independent of whether the synthetic images themselves fully cover the real data manifold. & 3.00 & 39 \\
\addlinespace[3pt]
18 & In black-box data-free KD, how can a student simultaneously (i) generate synthetic images that cover the teacher's full input distribution, and (ii) \textcolor{red}{\sout{recover}} a soft distillation signal \textcolor{red}{\sout{approximating}} the teacher's class probabilities, both without any real training data and with only hard-label access to the teacher? & In black-box data-free KD, how can a student simultaneously (i) generate synthetic images that cover the teacher's full input distribution, and (ii) \textcolor{blue}{\uwave{obtain}} a \textcolor{blue}{\uwave{surrogate}} soft distillation signal \textcolor{blue}{\uwave{in}} \textcolor{blue}{\uwave{place}} \textcolor{blue}{\uwave{of}} the teacher's class probabilities, both without any real training data and with only hard-label access to the teacher? & 3.50 & 40 \\
\addlinespace[3pt]
\multicolumn{5}{l}{\textbf{2.5 Distributional Fidelity and Synthetic Data $\rightarrow$ 2.5.1 The Fidelity Deficit}}\\[2pt]
19 & A surrogate composed entirely of easy-to-classify examples may cover the input space densely yet fail to represent the corpus distribution over the learning trajectories that matter: \textcolor{red}{\sout{and}} representative examples that capture the distributional core, boundary-adjacent examples that refine decisions. Coverage is therefore a necessary condition for distributional fidelity, not a separate objective; \textcolor{red}{\sout{and}} for dataset distillation (Chapter \#), fidelity to the full corpus distribution is the objective itself. & A surrogate composed entirely of easy-to-classify examples may cover the input space densely yet fail to represent the corpus distribution over the learning trajectories that matter, \textcolor{blue}{\uwave{which}} \textcolor{blue}{\uwave{requires}} \textcolor{blue}{\uwave{both}} representative examples that capture the distributional core \textcolor{blue}{\uwave{and}} boundary-adjacent examples that refine decisions. Coverage is therefore a necessary condition for distributional fidelity, not a separate objective. For dataset distillation (Chapter \#), fidelity to the full corpus distribution is the objective itself. & 2.75 & 41 \\
\addlinespace[3pt]
\multicolumn{5}{l}{\textbf{2.5 Distributional Fidelity and Synthetic Data $\rightarrow$ 2.5.2 Generation Mechanisms}}\\[2pt]
20 & WGAN-GP is the generative backbone for Chapters \textcolor{red}{\sout{\#}} \textcolor{red}{\sout{and}} \#, chosen for its stability and natural alignment with distributional fidelity objectives. & WGAN-GP is the generative backbone for Chapter \#, chosen for its stability and natural alignment with distributional fidelity objectives. & 1.25 & 42 \\
\addlinespace[3pt]
21 & The remaining challenge is distributional fidelity: even \textcolor{red}{\sout{an}} well-fine-tuned LLM may oversample frequent patterns and neglect rare but informative examples, producing a fluent yet \textcolor{red}{\sout{distributional-unfaithful}} surrogate. & The remaining challenge is distributional fidelity: even \textcolor{blue}{\uwave{a}} well-fine-tuned LLM may oversample frequent patterns and neglect rare but informative examples, producing a fluent yet \textcolor{blue}{\uwave{distributionally}} \textcolor{blue}{\uwave{unfaithful}} surrogate. & 2.25 & 42 \\
\addlinespace[3pt]
\multicolumn{5}{l}{\textbf{2.5 Distributional Fidelity and Synthetic Data $\rightarrow$ 2.5.3 Diagnosing Coverage Failures}}\\[2pt]
22 & Measuring coverage requires separating fidelity (quality of individual examples) from diversity (breadth of distributional support); standard metrics conflate the two. & Measuring coverage requires separating \textcolor{blue}{\uwave{sample}} fidelity (quality of individual examples) from diversity (breadth of distributional support); standard metrics conflate the two. & 1.00 & 43 \\
\addlinespace[3pt]
23 & Density measures fidelity: the average distance from each generated sample to its nearest real neighbor in embedding space (lower is better; generated samples are close to the real distribution). & Density measures \textcolor{blue}{\uwave{sample}} fidelity: the average distance from each generated sample to its nearest real neighbor in embedding space (lower is better; generated samples are close to the real distribution). & 1.00 & 43 \\
\addlinespace[3pt]
\multicolumn{5}{l}{\textbf{2.5 Distributional Fidelity and Synthetic Data $\rightarrow$ 2.5.4 Active Fidelity Recovery}}\\[2pt]
24 & High teacher confidence indicates that the sample is in a region the teacher \textcolor{red}{\sout{knows}} well; \textcolor{red}{\sout{rejection}} sampling \textcolor{red}{\sout{on}} \textcolor{red}{\sout{low-confidence}} \textcolor{red}{\sout{samples}} \textcolor{red}{\sout{pushes}} the \textcolor{red}{\sout{generator}} \textcolor{red}{\sout{toward}} \textcolor{red}{\sout{unexplored}} \textcolor{red}{\sout{regions}} \textcolor{red}{\sout{where}} the teacher's \textcolor{red}{\sout{decision}} \textcolor{red}{\sout{boundary}} \textcolor{red}{\sout{is}} \textcolor{red}{\sout{uncertain.}} & High teacher confidence indicates that the sample is in a region the teacher \textcolor{blue}{\uwave{recognizes}} well. \textcolor{blue}{\uwave{Admitting}} \textcolor{blue}{\uwave{such}} samples \textcolor{blue}{\uwave{as}} \textcolor{blue}{\uwave{additional}} \textcolor{blue}{\uwave{real}} \textcolor{blue}{\uwave{images}} \textcolor{blue}{\uwave{gives}} the \textcolor{blue}{\uwave{discriminator}} \textcolor{blue}{\uwave{signal}} \textcolor{blue}{\uwave{from}} \textcolor{blue}{\uwave{outside}} \textcolor{blue}{\uwave{the}} \textcolor{blue}{\uwave{few-shot}} \textcolor{blue}{\uwave{set,}} \textcolor{blue}{\uwave{with}} \textcolor{blue}{\uwave{confidence}} \textcolor{blue}{\uwave{serving}} \textcolor{blue}{\uwave{as}} \textcolor{blue}{\uwave{a}} \textcolor{blue}{\uwave{practical}} \textcolor{blue}{\uwave{proxy}} \textcolor{blue}{\uwave{for}} the teacher's \textcolor{blue}{\uwave{input}} \textcolor{blue}{\uwave{distribution.}} & 20.25 & 44 \\
\addlinespace[3pt]
\multicolumn{5}{l}{\textbf{2.9 Summary and Research Gaps}}\\[2pt]
25 & Coverage failures manifest differently, mode collapse in GAN-based synthesis and coverage gaps in clustering-based corpus selection, but stem from the same root cause: the optimization objectives used to produce the artifact do not directly \textcolor{red}{\sout{penalize}} distributional \textcolor{red}{\sout{divergence.}} The three contributions in Chapters \#--\# \textcolor{red}{\sout{each}} \textcolor{red}{\sout{make}} \textcolor{red}{\sout{direct}} \textcolor{red}{\sout{divergence}} \textcolor{red}{\sout{minimization}} \textcolor{red}{\sout{tractable}} in \textcolor{red}{\sout{its}} respective context. & Coverage failures manifest differently, mode collapse in GAN-based synthesis and coverage gaps in clustering-based corpus selection, but stem from the same root cause: the optimization objectives used to produce the artifact do not directly \textcolor{blue}{\uwave{promote}} distributional \textcolor{blue}{\uwave{coverage.}} The three contributions in Chapters \#--\# \textcolor{blue}{\uwave{address}} \textcolor{blue}{\uwave{this}} in \textcolor{blue}{\uwave{their}} respective contexts: \textcolor{blue}{\uwave{Chapters}} \textcolor{blue}{\uwave{\#}} \textcolor{blue}{\uwave{and}} \textcolor{blue}{\uwave{\#}} \textcolor{blue}{\uwave{actively}} \textcolor{blue}{\uwave{promote}} \textcolor{blue}{\uwave{distributional}} \textcolor{blue}{\uwave{coverage,}} \textcolor{blue}{\uwave{and}} \textcolor{blue}{\uwave{Chapter}} \textcolor{blue}{\uwave{\#}} \textcolor{blue}{\uwave{minimizes}} \textcolor{blue}{\uwave{a}} \textcolor{blue}{\uwave{transport}} \textcolor{blue}{\uwave{cost}} \textcolor{blue}{\uwave{to}} \textcolor{blue}{\uwave{the}} \textcolor{blue}{\uwave{source}} \textcolor{blue}{\uwave{distribution.}} & 21.00 & 58 \\
\addlinespace[3pt]
26 & In black-box few-shot KD, the student trains on a small image set and receives only \textcolor{red}{\sout{top-1}} \textcolor{red}{\sout{predictions}} from a black-box teacher. & In black-box few-shot KD, the student trains on a small \textcolor{blue}{\uwave{labeled}} image set and receives only \textcolor{blue}{\uwave{class}} \textcolor{blue}{\uwave{probabilities}} from a black-box teacher. & 3.00 & 59 \\
\addlinespace[3pt]
27 & The binding constraint is diversity: synthetic samples must venture beyond the few-shot set's inherent modes to provide \textcolor{red}{\sout{decision-boundary}} coverage. & The binding constraint is diversity: synthetic samples must venture beyond the few-shot set's inherent modes to provide \textcolor{blue}{\uwave{broader}} coverage \textcolor{blue}{\uwave{of}} \textcolor{blue}{\uwave{the}} \textcolor{blue}{\uwave{teacher's}} \textcolor{blue}{\uwave{input}} \textcolor{blue}{\uwave{space.}} & 4.50 & 59 \\
\addlinespace[3pt]
28 & Chapter \# (DIPKD) tackles both simultaneously via a three-phase pipeline: (1) Synthesis of diverse image priors using multi-scale noise; (2) Contrast via a primer student (white-box mediator) that enforces distinction via contrastive learning; (3) Distillation where the primer student \textcolor{red}{\sout{recovers}} soft \textcolor{red}{\sout{probabilistic}} \textcolor{red}{\sout{signals,}} \textcolor{red}{\sout{enabling}} \textcolor{red}{\sout{dark-knowledge}} \textcolor{red}{\sout{transfer}} \textcolor{red}{\sout{despite}} hard-label teacher access. & Chapter \# (DIPKD) tackles both simultaneously via a three-phase pipeline: (1) Synthesis of diverse image priors using multi-scale noise; (2) Contrast via a primer student (white-box mediator) that enforces distinction via contrastive learning; (3) Distillation where the primer student \textcolor{blue}{\uwave{supplies}} \textcolor{blue}{\uwave{surrogate}} soft \textcolor{blue}{\uwave{supervision,}} \textcolor{blue}{\uwave{standing}} \textcolor{blue}{\uwave{in}} \textcolor{blue}{\uwave{for}} \textcolor{blue}{\uwave{the}} \textcolor{blue}{\uwave{dark}} \textcolor{blue}{\uwave{knowledge}} \textcolor{blue}{\uwave{that}} hard-label teacher access \textcolor{blue}{\uwave{withholds.}} & 9.50 & 59 \\
\addlinespace[3pt]
\midrule \multicolumn{5}{l}{\textbf{\large Front matter}}\\
\midrule
\multicolumn{5}{l}{\textbf{3.1 Introduction}}\\[2pt]
29 & Then, during the WGAN training loop, we select high-confidence images from the latest iteration on-the-fly and introduce them in the role of additional real images, allowing our WGAN to learn and generate diverse synthetic images \textcolor{red}{\sout{close}} \textcolor{red}{\sout{to}} the teacher's \textcolor{red}{\sout{training}} \textcolor{red}{\sout{images.}} & Then, during the WGAN training loop, we select high-confidence images from the latest iteration on-the-fly and introduce them in the role of additional real images, allowing our WGAN to learn and generate diverse synthetic images \textcolor{blue}{\uwave{that}} \textcolor{blue}{\uwave{resemble}} \textcolor{blue}{\uwave{what}} the teacher \textcolor{blue}{\uwave{recognizes}} \textcolor{blue}{\uwave{confidently.}} & 3.50 & 62 \\
\addlinespace[3pt]
30 & Adaptive per-class confidence thresholds that calibrate selection to each class's \textcolor{red}{\sout{decision}} \textcolor{red}{\sout{boundary,}} preventing bias toward easy \textcolor{red}{\sout{cases.}} & Adaptive per-class confidence thresholds that calibrate selection to each class's \textcolor{blue}{\uwave{confidence}} \textcolor{blue}{\uwave{distribution,}} preventing bias toward easy \textcolor{blue}{\uwave{classes.}} & 3.00 & 63 \\
\addlinespace[3pt]
\multicolumn{5}{l}{\textbf{3.4 Conclusion}}\\[2pt]
31 & Through comprehensive evaluation, we have shown that our DivBFKD could use diversity to achieve \textcolor{red}{\sout{state-of-the-art}} \textcolor{red}{\sout{KD}} \textcolor{red}{\sout{performance}} among \textcolor{red}{\sout{other}} few-shot KD methods on multiple image datasets and architectures at different scales. & Through comprehensive evaluation, we have shown that our DivBFKD could use diversity to achieve \textcolor{blue}{\uwave{the}} \textcolor{blue}{\uwave{highest}} \textcolor{blue}{\uwave{accuracy}} among \textcolor{blue}{\uwave{the}} \textcolor{blue}{\uwave{compared}} few-shot KD methods \textcolor{blue}{\uwave{(BBKD}} \textcolor{blue}{\uwave{and}} \textcolor{blue}{\uwave{FS-BBT)}} on multiple image datasets and architectures at different scales. \textcolor{blue}{\uwave{As}} \textcolor{blue}{\uwave{a}} \textcolor{blue}{\uwave{limitation,}} \textcolor{blue}{\uwave{the}} \textcolor{blue}{\uwave{adversarial}} \textcolor{blue}{\uwave{training}} \textcolor{blue}{\uwave{of}} \textcolor{blue}{\uwave{the}} \textcolor{blue}{\uwave{WGAN}} \textcolor{blue}{\uwave{remains}} \textcolor{blue}{\uwave{sensitive}} \textcolor{blue}{\uwave{to}} \textcolor{blue}{\uwave{the}} \textcolor{blue}{\uwave{balance}} \textcolor{blue}{\uwave{between}} \textcolor{blue}{\uwave{generator}} \textcolor{blue}{\uwave{and}} \textcolor{blue}{\uwave{discriminator,}} \textcolor{blue}{\uwave{and}} \textcolor{blue}{\uwave{we}} \textcolor{blue}{\uwave{have}} \textcolor{blue}{\uwave{not}} \textcolor{blue}{\uwave{characterized}} \textcolor{blue}{\uwave{this}} \textcolor{blue}{\uwave{at}} \textcolor{blue}{\uwave{larger}} \textcolor{blue}{\uwave{resolutions.}} & 23.00 & 80 \\
\addlinespace[3pt]
\multicolumn{5}{l}{\textbf{4.1 Introduction}}\\[2pt]
32 & (1) Synthesis: crafts synthetic image priors with diverse patterns \textcolor{red}{\sout{and}} \textcolor{red}{\sout{semantics}} \textcolor{red}{\sout{of}} natural objects; & (1) Synthesis: crafts synthetic image priors with diverse \textcolor{blue}{\uwave{structural}} patterns \textcolor{blue}{\uwave{resembling}} natural objects; & 2.50 & 82 \\
\addlinespace[3pt]
33 & We introduce image priors, a novel class of synthetic images exhibiting diversity across hierarchical structures, nonlinearity, and \textcolor{red}{\sout{semantics.}} We propose the novel use of a primer student, enabling contrastive optimization and \textcolor{red}{\sout{extraction}} \textcolor{red}{\sout{of}} soft \textcolor{red}{\sout{probabilistic}} \textcolor{red}{\sout{signals}} \textcolor{red}{\sout{originally}} \textcolor{red}{\sout{restricted}} \textcolor{red}{\sout{in}} black-box \textcolor{red}{\sout{KD.}} Extensive analysis of our framework \textcolor{red}{\sout{demonstrates}} that it \textcolor{red}{\sout{effectively}} expands the \textcolor{red}{\sout{semantic}} coverage of synthetic domains, leading to \textcolor{red}{\sout{our}} state-of-the-art performance \textcolor{red}{\sout{across}} 12 benchmarks. & We introduce image priors, a novel class of synthetic images exhibiting diversity across hierarchical structures, nonlinearity, and \textcolor{blue}{\uwave{spatial}} \textcolor{blue}{\uwave{structure.}} We propose the novel use of a primer student, enabling contrastive optimization and \textcolor{blue}{\uwave{surrogate}} soft \textcolor{blue}{\uwave{supervision}} \textcolor{blue}{\uwave{where}} \textcolor{blue}{\uwave{the}} black-box \textcolor{blue}{\uwave{teacher}} \textcolor{blue}{\uwave{provides}} \textcolor{blue}{\uwave{only}} \textcolor{blue}{\uwave{hard}} \textcolor{blue}{\uwave{labels.}} Extensive analysis of our framework \textcolor{blue}{\uwave{shows}} that it expands the coverage of synthetic domains, leading to state-of-the-art performance \textcolor{blue}{\uwave{on}} \textcolor{blue}{\uwave{11}} \textcolor{blue}{\uwave{of}} 12 benchmarks. & 15.50 & 82 \\
\addlinespace[3pt]
\multicolumn{5}{l}{\textbf{4.2 Proposed Framework $\rightarrow$ 4.2.1 Problem setup}}\\[2pt]
34 & In the Synthesis phase, we design a pipeline to create image priors $\mathcal{X}_\mathcal{P}$ that simulate the hierarchical structures, nonlinearity, and \textcolor{red}{\sout{meaningful}} \textcolor{red}{\sout{semantics}} in natural scenes and objects. & In the Synthesis phase, we design a pipeline to create image priors $\mathcal{X}_\mathcal{P}$ that simulate the hierarchical structures, nonlinearity, and \textcolor{blue}{\uwave{spatial}} \textcolor{blue}{\uwave{structure}} in natural scenes and objects. & 2.00 & 84 \\
\addlinespace[3pt]
35 & Crucially, we extract soft knowledge from the primer student $S_0$ \textcolor{red}{\sout{to}} \textcolor{red}{\sout{restore}} the inter-class signals absent in previous BBDFKD methodologies. & Crucially, we extract soft knowledge from the primer student $S_0$ \textcolor{blue}{\uwave{as}} \textcolor{blue}{\uwave{surrogate}} \textcolor{blue}{\uwave{supervision}} \textcolor{blue}{\uwave{for}} the inter-class signals absent in previous BBDFKD methodologies. & 2.50 & 84 \\
\addlinespace[3pt]
\multicolumn{5}{l}{\textbf{4.2 Proposed Framework $\rightarrow$ 4.2.2 Synthesis: Reconstructing the knowledge domain}}\\[2pt]
36 & and (3) Semantic cutmixing: We apply cutmixing on $x_\text{nt}$ with a semantic mask $\hat{m}$ to construct the final images $x$, simulating diverse \textcolor{red}{\sout{semantics.}} & and (3) Semantic cutmixing: We apply cutmixing on $x_\text{nt}$ with a semantic mask $\hat{m}$ to construct the final images $x$, simulating diverse \textcolor{blue}{\uwave{spatial}} \textcolor{blue}{\uwave{structures.}} & 2.00 & 84 \\
\addlinespace[3pt]
37 & This pipeline serves to capture the structural \textcolor{red}{\sout{and}} \textcolor{red}{\sout{semantic}} breadth of the expert's original knowledge domain through enhanced diversity. & This pipeline serves to capture the structural breadth of the expert's original knowledge domain through enhanced diversity. & 1.25 & 85 \\
\addlinespace[3pt]
\multicolumn{5}{l}{\textbf{4.2 Proposed Framework $\rightarrow$ 4.2.2.3 Semantic cutmixing}}\\[2pt]
38 & To address this, we extend CutMix by improving the masking process to nonlinearly blend image pairs into \textcolor{red}{\sout{semantically}} coherent structures. & To address this, we extend CutMix by improving the masking process to nonlinearly blend image pairs into \textcolor{blue}{\uwave{irregular,}} \textcolor{blue}{\uwave{spatially}} coherent structures, \textcolor{blue}{\uwave{which}} \textcolor{blue}{\uwave{we}} \textcolor{blue}{\uwave{call}} \textcolor{blue}{\uwave{semantic}} \textcolor{blue}{\uwave{cutmix.}} \textcolor{blue}{\uwave{These}} \textcolor{blue}{\uwave{masks}} \textcolor{blue}{\uwave{are}} \textcolor{blue}{\uwave{structural:}} \textcolor{blue}{\uwave{they}} \textcolor{blue}{\uwave{are}} \textcolor{blue}{\uwave{not}} \textcolor{blue}{\uwave{object-level}} \textcolor{blue}{\uwave{segments,}} \textcolor{blue}{\uwave{and}} \textcolor{blue}{\uwave{a}} \textcolor{blue}{\uwave{teacher}} \textcolor{blue}{\uwave{label}} \textcolor{blue}{\uwave{on}} \textcolor{blue}{\uwave{a}} \textcolor{blue}{\uwave{blended}} \textcolor{blue}{\uwave{image}} \textcolor{blue}{\uwave{does}} \textcolor{blue}{\uwave{not}} \textcolor{blue}{\uwave{make}} \textcolor{blue}{\uwave{it}} \textcolor{blue}{\uwave{a}} \textcolor{blue}{\uwave{meaningful}} \textcolor{blue}{\uwave{class}} \textcolor{blue}{\uwave{instance.}} & 22.25 & 86 \\
\addlinespace[3pt]
39 & This approach \textcolor{red}{\sout{ensures}} the synthetic domain of $\mathcal{X}_\mathcal{P}$ \textcolor{red}{\sout{captures}} diverse \textcolor{red}{\sout{semantics}} \textcolor{red}{\sout{required}} for \textcolor{red}{\sout{effective}} black-box knowledge acquisition. & This approach \textcolor{blue}{\uwave{encourages}} the synthetic domain of $\mathcal{X}_\mathcal{P}$ \textcolor{blue}{\uwave{to}} \textcolor{blue}{\uwave{capture}} diverse \textcolor{blue}{\uwave{structural}} \textcolor{blue}{\uwave{patterns}} \textcolor{blue}{\uwave{useful}} for black-box knowledge acquisition. & 5.75 & 86 \\
\addlinespace[3pt]
\multicolumn{5}{l}{\textbf{4.2 Proposed Framework $\rightarrow$ 4.2.4 Distillation: Restoring informative signals}}\\[2pt]
40 & Our objective is to \textcolor{red}{\sout{restore}} the high-bandwidth signals typical of standard KD [cit.], , incorporating both hard and soft knowledge, which were limited or entirely absent in recent works [cit.]. Having privileged access to the primer network $S_0$, we form \textcolor{red}{\sout{an}} \textcolor{red}{\sout{optimal}} distillation objective by jointly matching the student's logits $S^\mathcal{Z}(x)$ with soft labels $S_{0}^\mathcal{Z}(x)$ and hard labels $\hat{y}_{T}^\text{hard}$: & Our objective is to \textcolor{blue}{\uwave{approximate}} the high-bandwidth signals typical of standard KD [cit.], , incorporating both hard and soft knowledge, which were limited or entirely absent in recent works [cit.]. Having privileged access to the primer network $S_0$, we form \textcolor{blue}{\uwave{a}} distillation objective by jointly matching the student's logits $S^\mathcal{Z}(x)$ with soft labels $S_{0}^\mathcal{Z}(x)$ and hard labels $\hat{y}_{T}^\text{hard}$: & 2.25 & 88 \\
\addlinespace[3pt]
\multicolumn{5}{l}{\textbf{4.3 Experiments $\rightarrow$ 4.3.4.1 Contribution of components}}\\[2pt]
41 & \textit{(new text)} & \textcolor{blue}{\uwave{In}} \textcolor{blue}{\uwave{our}} \textcolor{blue}{\uwave{component}} \textcolor{blue}{\uwave{ablation}} \textcolor{blue}{\uwave{study,}} \textcolor{blue}{\uwave{the}} \textcolor{blue}{\uwave{Contrast}} \textcolor{blue}{\uwave{and}} \textcolor{blue}{\uwave{Distillation}} \textcolor{blue}{\uwave{phases}} \textcolor{blue}{\uwave{add}} \textcolor{blue}{\uwave{little}} \textcolor{blue}{\uwave{on}} \textcolor{blue}{\uwave{simple}} \textcolor{blue}{\uwave{datasets}} \textcolor{blue}{\uwave{but}} \textcolor{blue}{\uwave{clearly}} \textcolor{blue}{\uwave{help}} \textcolor{blue}{\uwave{on}} \textcolor{blue}{\uwave{complex}} \textcolor{blue}{\uwave{ones,}} \textcolor{blue}{\uwave{so}} \textcolor{blue}{\uwave{the}} \textcolor{blue}{\uwave{extra}} \textcolor{blue}{\uwave{training}} \textcolor{blue}{\uwave{beyond}} \textcolor{blue}{\uwave{reusing}} \textcolor{blue}{\uwave{the}} \textcolor{blue}{\uwave{primer}} \textcolor{blue}{\uwave{as}} \textcolor{blue}{\uwave{the}} \textcolor{blue}{\uwave{final}} \textcolor{blue}{\uwave{student}} \textcolor{blue}{\uwave{is}} \textcolor{blue}{\uwave{most}} \textcolor{blue}{\uwave{justified}} \textcolor{blue}{\uwave{there.}} & 27.25 & 91 \\
\addlinespace[3pt]
\multicolumn{5}{l}{\textbf{4.4 Conclusion}}\\[2pt]
42 & We achieve state-of-the-art performance \textcolor{red}{\sout{across}} 12 benchmarks, with notable improvements in complex and domain-specific settings. & We achieve state-of-the-art performance \textcolor{blue}{\uwave{on}} \textcolor{blue}{\uwave{11}} \textcolor{blue}{\uwave{of}} 12 benchmarks, with notable improvements in complex and domain-specific settings. & 1.50 & 96 \\
\addlinespace[3pt]
\multicolumn{5}{l}{\textbf{5.4 Experiments $\rightarrow$ 5.4.1 Setup}}\\[2pt]
43 & We evaluate on two linguistic tasks: AG News [cit.], IMDb [cit.], SST-2 for classification; MNLI-m, QQP, and QNLI for NLI. & We evaluate on two linguistic tasks: AG News [cit.], IMDb [cit.], SST-2 for classification; MNLI-m, QQP, and QNLI for NLI \textcolor{blue}{\uwave{(QQP}} \textcolor{blue}{\uwave{is}} \textcolor{blue}{\uwave{a}} \textcolor{blue}{\uwave{sentence-pair}} \textcolor{blue}{\uwave{duplicate-detection}} \textcolor{blue}{\uwave{task).}} & 4.50 & 111 \\
\addlinespace[3pt]
\multicolumn{5}{l}{\textbf{5.5 Limitations and Ethical Considerations}}\\[2pt]
44 & The $L$-Lipschitz assumption on $\ell$ (Theorem \textcolor{red}{\sout{1)}} holds approximately when $\phi$ is task-aligned and embeddings are normalized, but remains an approximation for cross-entropy over discrete text. & The $L$-Lipschitz assumption on $\ell$ (Theorem \textcolor{blue}{\uwave{\#)}} holds approximately when $\phi$ is task-aligned and embeddings are normalized, but remains an approximation for cross-entropy over discrete text. & 1.00 & 116 \\
\addlinespace[3pt]
\midrule \multicolumn{5}{l}{\textbf{\large 6 Conclusion}}\\
\midrule
\multicolumn{5}{l}{\textbf{(chapter opening)}}\\[2pt]
45 & Each contribution treats fidelity not as a property to hope for at sufficient scale, but as an objective to enforce by design: through adaptive confidence filtering, through structurally diverse synthesis priors, and through trajectory-aware influence weighting coupled with \textcolor{red}{\sout{globally}} \textcolor{red}{\sout{optimal}} \textcolor{red}{\sout{transport.}} & Each contribution treats fidelity not as a property to hope for at sufficient scale, but as an objective to enforce by design: through adaptive confidence filtering, through structurally diverse synthesis priors, and through trajectory-aware influence weighting coupled with \textcolor{blue}{\uwave{a}} \textcolor{blue}{\uwave{set-level}} \textcolor{blue}{\uwave{optimal-transport}} \textcolor{blue}{\uwave{objective}} \textcolor{blue}{\uwave{that}} \textcolor{blue}{\uwave{avoids}} \textcolor{blue}{\uwave{redundant}} \textcolor{blue}{\uwave{selection.}} & 6.50 & 117 \\
\addlinespace[3pt]
\multicolumn{5}{l}{\textbf{6.1 Summary of Contributions}}\\[2pt]
46 & DivBFKD (Chapter \#) addresses the few-shot black-box KD problem, where the student has access to only a small \textcolor{red}{\sout{unlabeled}} image set and may query the teacher solely for \textcolor{red}{\sout{top-1}} \textcolor{red}{\sout{predictions.}} & DivBFKD (Chapter \#) addresses the few-shot black-box KD problem, where the student has access to only a small \textcolor{blue}{\uwave{labeled}} image set and may query the teacher solely for \textcolor{blue}{\uwave{its}} \textcolor{blue}{\uwave{predicted}} \textcolor{blue}{\uwave{class}} \textcolor{blue}{\uwave{probabilities.}} & 4.25 & 117 \\
\addlinespace[3pt]
47 & Adaptive per-class \textcolor{red}{\sout{confidence}} thresholds \textcolor{red}{\sout{prevent}} the teacher's class-wise \textcolor{red}{\sout{accuracy}} bias \textcolor{red}{\sout{from}} \textcolor{red}{\sout{skewing}} \textcolor{red}{\sout{which}} \textcolor{red}{\sout{images}} \textcolor{red}{\sout{are}} \textcolor{red}{\sout{admitted,}} \textcolor{red}{\sout{ensuring}} \textcolor{red}{\sout{that}} \textcolor{red}{\sout{the}} \textcolor{red}{\sout{selected}} \textcolor{red}{\sout{exemplars}} \textcolor{red}{\sout{represent}} \textcolor{red}{\sout{the}} \textcolor{red}{\sout{teacher's}} \textcolor{red}{\sout{decision}} \textcolor{red}{\sout{boundaries}} rather than \textcolor{red}{\sout{its}} \textcolor{red}{\sout{easiest}} \textcolor{red}{\sout{cases.}} \textcolor{red}{\sout{An}} \textcolor{red}{\sout{ablation}} \textcolor{red}{\sout{study}} \textcolor{red}{\sout{confirms}} \textcolor{red}{\sout{that}} \textcolor{red}{\sout{removing}} \textcolor{red}{\sout{adaptive}} \textcolor{red}{\sout{thresholds---reverting}} \textcolor{red}{\sout{to}} \textcolor{red}{\sout{a}} global \textcolor{red}{\sout{cutoff---degrades}} \textcolor{red}{\sout{performance,}} \textcolor{red}{\sout{validating}} \textcolor{red}{\sout{that}} \textcolor{red}{\sout{class-wise}} \textcolor{red}{\sout{calibration}} \textcolor{red}{\sout{of}} \textcolor{red}{\sout{the}} \textcolor{red}{\sout{filtering}} \textcolor{red}{\sout{criterion}} \textcolor{red}{\sout{is}} \textcolor{red}{\sout{load-bearing,}} \textcolor{red}{\sout{not}} \textcolor{red}{\sout{incidental.}} DivBFKD achieves \textcolor{red}{\sout{state-of-the-art}} \textcolor{red}{\sout{performance}} among few-shot black-box KD methods across \textcolor{red}{\sout{multiple}} image classification datasets and architectures. & Adaptive per-class thresholds \textcolor{blue}{\uwave{keep}} \textcolor{blue}{\uwave{the}} \textcolor{blue}{\uwave{easiest}} \textcolor{blue}{\uwave{classes}} \textcolor{blue}{\uwave{from}} \textcolor{blue}{\uwave{dominating,}} \textcolor{blue}{\uwave{so}} \textcolor{blue}{\uwave{each}} \textcolor{blue}{\uwave{class}} \textcolor{blue}{\uwave{contributes}} \textcolor{blue}{\uwave{reliably}} \textcolor{blue}{\uwave{labeled}} \textcolor{blue}{\uwave{exemplars}} \textcolor{blue}{\uwave{from}} \textcolor{blue}{\uwave{well}} \textcolor{blue}{\uwave{inside}} \textcolor{blue}{\uwave{the}} \textcolor{blue}{\uwave{teacher's}} \textcolor{blue}{\uwave{class}} \textcolor{blue}{\uwave{regions}} \textcolor{blue}{\uwave{to}} \textcolor{blue}{\uwave{anchor}} \textcolor{blue}{\uwave{the}} \textcolor{blue}{\uwave{generator.}} \textcolor{blue}{\uwave{An}} \textcolor{blue}{\uwave{ablation}} \textcolor{blue}{\uwave{study}} \textcolor{blue}{\uwave{shows}} \textcolor{blue}{\uwave{that}} \textcolor{blue}{\uwave{the}} \textcolor{blue}{\uwave{adaptive-threshold}} \textcolor{blue}{\uwave{scheme}} \textcolor{blue}{\uwave{improves}} \textcolor{blue}{\uwave{student}} \textcolor{blue}{\uwave{accuracy}} \textcolor{blue}{\uwave{over}} \textcolor{blue}{\uwave{a}} \textcolor{blue}{\uwave{plain}} \textcolor{blue}{\uwave{WGAN,}} \textcolor{blue}{\uwave{and}} the teacher's class-wise \textcolor{blue}{\uwave{confidence}} bias \textcolor{blue}{\uwave{motivates}} \textcolor{blue}{\uwave{per-class}} rather than global \textcolor{blue}{\uwave{thresholds.}} DivBFKD achieves \textcolor{blue}{\uwave{the}} \textcolor{blue}{\uwave{highest}} \textcolor{blue}{\uwave{accuracy}} among \textcolor{blue}{\uwave{the}} \textcolor{blue}{\uwave{compared}} few-shot black-box KD methods \textcolor{blue}{\uwave{(BBKD}} \textcolor{blue}{\uwave{and}} \textcolor{blue}{\uwave{FS-BBT)}} across \textcolor{blue}{\uwave{the}} image classification datasets and architectures \textcolor{blue}{\uwave{evaluated.}} & 46.00 & 118 \\
\addlinespace[3pt]
48 & The fidelity problem becomes entirely generative---the student must be trained on a synthetic distribution that covers the teacher's decision \textcolor{red}{\sout{surface}} without any real-image anchor. & The fidelity problem becomes entirely generative---the student must be trained on a synthetic distribution that \textcolor{blue}{\uwave{aims}} \textcolor{blue}{\uwave{to}} cover \textcolor{blue}{\uwave{the}} \textcolor{blue}{\uwave{inputs}} \textcolor{blue}{\uwave{relevant}} \textcolor{blue}{\uwave{to}} the teacher's decisions, without any real-image anchor. & 4.75 & 118 \\
\addlinespace[3pt]
49 & The Synthesis phase constructs hierarchical multi-scale image priors from uniform noise with \textcolor{red}{\sout{geometric}} \textcolor{red}{\sout{transforms}} and semantic cutmixing, producing structural diversity before any teacher query. The Contrast phase passes these priors through a primer student backbone and applies \textcolor{red}{\sout{a}} contrastive objective that optimizes the synthetic dataset directly, pushing \textcolor{red}{\sout{candidates}} \textcolor{red}{\sout{apart}} in embedding space to prevent mode collapse. The Distillation phase then \textcolor{red}{\sout{recovers}} soft \textcolor{red}{\sout{probabilistic}} \textcolor{red}{\sout{labels}} from the primer \textcolor{red}{\sout{student---approximating}} the dark knowledge the black-box teacher \textcolor{red}{\sout{withholds---and}} uses them alongside hard teacher labels to train the final student. DIPKD achieves state-of-the-art performance \textcolor{red}{\sout{across}} \textcolor{red}{\sout{twelve}} benchmarks, with \textcolor{red}{\sout{the}} \textcolor{red}{\sout{largest}} gains \textcolor{red}{\sout{in}} \textcolor{red}{\sout{fine-grained}} and domain-specific \textcolor{red}{\sout{settings}} where synthetic diversity is most constrained. & The Synthesis phase constructs hierarchical multi-scale image priors from uniform noise with \textcolor{blue}{\uwave{nonlinear}} \textcolor{blue}{\uwave{filters}} and semantic cutmixing, \textcolor{blue}{\uwave{whose}} \textcolor{blue}{\uwave{masks}} \textcolor{blue}{\uwave{are}} \textcolor{blue}{\uwave{irregular}} \textcolor{blue}{\uwave{regions}} \textcolor{blue}{\uwave{rather}} \textcolor{blue}{\uwave{than}} \textcolor{blue}{\uwave{object}} \textcolor{blue}{\uwave{segments,}} producing structural diversity before any teacher query. The Contrast phase passes these priors through a primer student backbone and applies \textcolor{blue}{\uwave{an}} \textcolor{blue}{\uwave{instance-level}} contrastive objective that optimizes the synthetic dataset directly, pushing \textcolor{blue}{\uwave{each}} \textcolor{blue}{\uwave{synthetic}} \textcolor{blue}{\uwave{image}} \textcolor{blue}{\uwave{away}} \textcolor{blue}{\uwave{from}} \textcolor{blue}{\uwave{all}} \textcolor{blue}{\uwave{others}} in embedding space to prevent mode collapse. The Distillation phase then \textcolor{blue}{\uwave{takes}} soft \textcolor{blue}{\uwave{outputs}} from the primer \textcolor{blue}{\uwave{student}} \textcolor{blue}{\uwave{as}} \textcolor{blue}{\uwave{surrogate}} \textcolor{blue}{\uwave{soft}} \textcolor{blue}{\uwave{supervision,}} \textcolor{blue}{\uwave{in}} \textcolor{blue}{\uwave{place}} \textcolor{blue}{\uwave{of}} the dark knowledge the black-box teacher \textcolor{blue}{\uwave{withholds,}} \textcolor{blue}{\uwave{and}} uses them alongside hard teacher labels to train the final student. DIPKD achieves state-of-the-art performance \textcolor{blue}{\uwave{on}} \textcolor{blue}{\uwave{11}} \textcolor{blue}{\uwave{of}} \textcolor{blue}{\uwave{the}} \textcolor{blue}{\uwave{12}} benchmarks \textcolor{blue}{\uwave{(all}} \textcolor{blue}{\uwave{except}} \textcolor{blue}{\uwave{Tiny-ImageNet),}} with \textcolor{blue}{\uwave{substantial}} gains \textcolor{blue}{\uwave{on}} \textcolor{blue}{\uwave{several}} \textcolor{blue}{\uwave{complex}} and domain-specific \textcolor{blue}{\uwave{datasets}} where synthetic diversity is most constrained. & 36.50 & 118 \\
\addlinespace[3pt]
50 & The goal is to replace a large labeled text corpus with a compact surrogate that \textcolor{red}{\sout{preserves}} downstream task performance. & The goal is to replace a large labeled text corpus with a compact surrogate that \textcolor{blue}{\uwave{approximates}} downstream task performance. & 1.00 & 119 \\
\addlinespace[3pt]
51 & Trajectory-aware influence scoring computes importance weights from reciprocal gradient norms accumulated across the training trajectory, correcting the hard-sample bias \textcolor{red}{\sout{that}} \textcolor{red}{\sout{convergence-point}} scoring \textcolor{red}{\sout{introduces}} \textcolor{red}{\sout{by}} assigning inflated importance to noisy or boundary \textcolor{red}{\sout{examples.}} & Trajectory-aware influence scoring computes importance weights from reciprocal gradient norms accumulated across the training trajectory, correcting the hard-sample bias, \textcolor{blue}{\uwave{whereby}} scoring \textcolor{blue}{\uwave{at}} \textcolor{blue}{\uwave{the}} \textcolor{blue}{\uwave{convergence}} \textcolor{blue}{\uwave{point}} \textcolor{blue}{\uwave{tends}} \textcolor{blue}{\uwave{to}} assign inflated importance to \textcolor{blue}{\uwave{hard}} \textcolor{blue}{\uwave{examples}} \textcolor{blue}{\uwave{such}} \textcolor{blue}{\uwave{as}} noisy or boundary \textcolor{blue}{\uwave{ones.}} & 8.50 & 119 \\
\addlinespace[3pt]
52 & The importance-weighted OT objective combines task alignment and distributional coverage as complementary objectives: importance alone \textcolor{red}{\sout{directs}} \textcolor{red}{\sout{the}} \textcolor{red}{\sout{budget}} \textcolor{red}{\sout{but}} \textcolor{red}{\sout{ignores}} \textcolor{red}{\sout{global}} \textcolor{red}{\sout{structure;}} \textcolor{red}{\sout{coverage}} \textcolor{red}{\sout{alone}} \textcolor{red}{\sout{ensures}} \textcolor{red}{\sout{spread}} \textcolor{red}{\sout{but}} \textcolor{red}{\sout{ignores}} \textcolor{red}{\sout{contribution;}} \textcolor{red}{\sout{together}} \textcolor{red}{\sout{they}} \textcolor{red}{\sout{close}} \textcolor{red}{\sout{failure}} \textcolor{red}{\sout{modes}} \textcolor{red}{\sout{that}} \textcolor{red}{\sout{neither}} \textcolor{red}{\sout{addresses}} \textcolor{red}{\sout{individually.}} & The importance-weighted OT objective combines task alignment and distributional coverage as complementary objectives, \textcolor{blue}{\uwave{and}} \textcolor{blue}{\uwave{in}} \textcolor{blue}{\uwave{our}} \textcolor{blue}{\uwave{ablation,}} \textcolor{blue}{\uwave{adding}} importance \textcolor{blue}{\uwave{weights}} \textcolor{blue}{\uwave{to}} \textcolor{blue}{\uwave{the}} \textcolor{blue}{\uwave{OT}} \textcolor{blue}{\uwave{objective}} \textcolor{blue}{\uwave{improves}} \textcolor{blue}{\uwave{over}} \textcolor{blue}{\uwave{OT}} alone. & 22.75 & 119 \\
\addlinespace[3pt]
53 & Across all three contributions, the fidelity-first design stance \textcolor{red}{\sout{enables}} deployment in regimes that scale-dependent assumptions foreclose: edge and embedded devices for DivBFKD and DIPKD, and low-resource, privacy-constrained, and carbon-sensitive language model development for TAKE. & \textcolor{blue}{\uwave{Table}} \textcolor{blue}{\uwave{\#}} \textcolor{blue}{\uwave{compares}} \textcolor{blue}{\uwave{the}} \textcolor{blue}{\uwave{three}} \textcolor{blue}{\uwave{contributions}} \textcolor{blue}{\uwave{by}} \textcolor{blue}{\uwave{assumptions,}} \textcolor{blue}{\uwave{information,}} \textcolor{blue}{\uwave{mechanism,}} \textcolor{blue}{\uwave{key}} \textcolor{blue}{\uwave{result,}} \textcolor{blue}{\uwave{and}} \textcolor{blue}{\uwave{limitation.}} \textcolor{blue}{\uwave{Comparison}} \textcolor{blue}{\uwave{of}} \textcolor{blue}{\uwave{the}} \textcolor{blue}{\uwave{three}} \textcolor{blue}{\uwave{thesis}} \textcolor{blue}{\uwave{contributions.}} \textcolor{blue}{\uwave{Assumptions}} \textcolor{blue}{\uwave{\&}} \textcolor{blue}{\uwave{A}} \textcolor{blue}{\uwave{small}} \textcolor{blue}{\uwave{labeled}} \textcolor{blue}{\uwave{image}} \textcolor{blue}{\uwave{set.}} \textcolor{blue}{\uwave{A}} \textcolor{blue}{\uwave{black-box}} \textcolor{blue}{\uwave{teacher.}} \textcolor{blue}{\uwave{\&}} \textcolor{blue}{\uwave{No}} \textcolor{blue}{\uwave{real}} \textcolor{blue}{\uwave{data.}} \textcolor{blue}{\uwave{A}} \textcolor{blue}{\uwave{black-box}} \textcolor{blue}{\uwave{teacher.}} \textcolor{blue}{\uwave{Information}} \textcolor{blue}{\uwave{\&}} \textcolor{blue}{\uwave{Class}} \textcolor{blue}{\uwave{probabilities}} \textcolor{blue}{\uwave{for}} \textcolor{blue}{\uwave{queried}} \textcolor{blue}{\uwave{images.}} \textcolor{blue}{\uwave{\&}} \textcolor{blue}{\uwave{Top-1}} \textcolor{blue}{\uwave{label}} \textcolor{blue}{\uwave{for}} \textcolor{blue}{\uwave{queried}} \textcolor{blue}{\uwave{images.}} \textcolor{blue}{\uwave{Mechanism}} \textcolor{blue}{\uwave{\&}} \textcolor{blue}{\uwave{WGAN}} \textcolor{blue}{\uwave{with}} \textcolor{blue}{\uwave{high-confidence}} \textcolor{blue}{\uwave{images}} \textcolor{blue}{\uwave{as}} \textcolor{blue}{\uwave{extra}} \textcolor{blue}{\uwave{real}} \textcolor{blue}{\uwave{images.}} \textcolor{blue}{\uwave{Adaptive}} \textcolor{blue}{\uwave{per-class}} \textcolor{blue}{\uwave{thresholds.}} \textcolor{blue}{\uwave{\&}} \textcolor{blue}{\uwave{Image}} \textcolor{blue}{\uwave{priors}} \textcolor{blue}{\uwave{built}} \textcolor{blue}{\uwave{from}} \textcolor{blue}{\uwave{noise.}} \textcolor{blue}{\uwave{Instance-level}} \textcolor{blue}{\uwave{contrastive}} \textcolor{blue}{\uwave{learning.}} \textcolor{blue}{\uwave{A}} \textcolor{blue}{\uwave{primer}} \textcolor{blue}{\uwave{student}} \textcolor{blue}{\uwave{for}} \textcolor{blue}{\uwave{surrogate}} \textcolor{blue}{\uwave{soft}} \textcolor{blue}{\uwave{labels.}} \textcolor{blue}{\uwave{Key}} \textcolor{blue}{\uwave{result}} \textcolor{blue}{\uwave{\&}} \textcolor{blue}{\uwave{Highest}} \textcolor{blue}{\uwave{accuracy}} \textcolor{blue}{\uwave{among}} \textcolor{blue}{\uwave{the}} \textcolor{blue}{\uwave{compared}} \textcolor{blue}{\uwave{few-shot}} \textcolor{blue}{\uwave{methods.}} \textcolor{blue}{\uwave{A}} \textcolor{blue}{\uwave{17\%}} \textcolor{blue}{\uwave{gain}} \textcolor{blue}{\uwave{over}} \textcolor{blue}{\uwave{standard}} \textcolor{blue}{\uwave{KD}} \textcolor{blue}{\uwave{on}} \textcolor{blue}{\uwave{CIFAR-10.}} \textcolor{blue}{\uwave{\&}} \textcolor{blue}{\uwave{State}} \textcolor{blue}{\uwave{of}} \textcolor{blue}{\uwave{the}} \textcolor{blue}{\uwave{art}} \textcolor{blue}{\uwave{on}} \textcolor{blue}{\uwave{11}} \textcolor{blue}{\uwave{of}} \textcolor{blue}{\uwave{12}} \textcolor{blue}{\uwave{benchmarks.}} \textcolor{blue}{\uwave{Tiny-ImageNet}} \textcolor{blue}{\uwave{is}} \textcolor{blue}{\uwave{the}} \textcolor{blue}{\uwave{exception.}} \textcolor{blue}{\uwave{Limitation}} \textcolor{blue}{\uwave{\&}} \textcolor{blue}{\uwave{Confidence}} \textcolor{blue}{\uwave{is}} \textcolor{blue}{\uwave{a}} \textcolor{blue}{\uwave{proxy}} \textcolor{blue}{\uwave{for}} \textcolor{blue}{\uwave{coverage.}} \textcolor{blue}{\uwave{Training}} \textcolor{blue}{\uwave{queries}} \textcolor{blue}{\uwave{the}} \textcolor{blue}{\uwave{teacher}} \textcolor{blue}{\uwave{more.}} \textcolor{blue}{\uwave{GAN}} \textcolor{blue}{\uwave{training}} \textcolor{blue}{\uwave{can}} \textcolor{blue}{\uwave{be}} \textcolor{blue}{\uwave{unstable.}} \textcolor{blue}{\uwave{\&}} \textcolor{blue}{\uwave{Priors}} \textcolor{blue}{\uwave{are}} \textcolor{blue}{\uwave{structurally}} \textcolor{blue}{\uwave{diverse}} \textcolor{blue}{\uwave{but}} \textcolor{blue}{\uwave{not}} \textcolor{blue}{\uwave{semantic.}} \textcolor{blue}{\uwave{The}} \textcolor{blue}{\uwave{pipeline}} \textcolor{blue}{\uwave{needs}} \textcolor{blue}{\uwave{extra}} \textcolor{blue}{\uwave{queries}} \textcolor{blue}{\uwave{and}} \textcolor{blue}{\uwave{training.}} Across all three contributions, the fidelity-first design stance \textcolor{blue}{\uwave{opens}} \textcolor{blue}{\uwave{the}} \textcolor{blue}{\uwave{potential}} \textcolor{blue}{\uwave{for}} deployment in regimes that scale-dependent assumptions foreclose: edge and embedded devices for DivBFKD and DIPKD, and low-resource, privacy-constrained, and carbon-sensitive language model development for TAKE. & 110.25 & 119 \\
\addlinespace[3pt]
\multicolumn{5}{l}{\textbf{6.2 Limitations and Future Directions $\rightarrow$ 6.2.1 Technical Limitations}}\\[2pt]
54 & More fundamentally, high-confidence filtering \textcolor{red}{\sout{guarantees}} teacher \textcolor{red}{\sout{agreement}} but not task-relevant diversity: a set of confident images may still be clustered in feature space if the teacher is confidently right for consistent visual reasons. & More fundamentally, high-confidence filtering \textcolor{blue}{\uwave{favors}} \textcolor{blue}{\uwave{images}} \textcolor{blue}{\uwave{on}} \textcolor{blue}{\uwave{which}} \textcolor{blue}{\uwave{the}} teacher \textcolor{blue}{\uwave{is}} \textcolor{blue}{\uwave{confident}} but \textcolor{blue}{\uwave{does}} not \textcolor{blue}{\uwave{ensure}} task-relevant diversity: a set of confident images may still be clustered in feature space if the teacher is confidently right for consistent visual reasons. & 5.25 & 121 \\
\addlinespace[3pt]
55 & DIPKD's Synthesis phase generates priors from hierarchical multi-scale uniform noise with \textcolor{red}{\sout{geometric}} \textcolor{red}{\sout{transforms}} and semantic cutmixing. This procedure yields structural diversity but is semantically agnostic: the priors may not cover the natural image manifold in \textcolor{red}{\sout{fine-grained}} or domain-specific tasks where intra-class variation is subtle. & DIPKD's Synthesis phase generates priors from hierarchical multi-scale uniform noise with \textcolor{blue}{\uwave{nonlinear}} \textcolor{blue}{\uwave{filters}} and semantic cutmixing. This procedure yields structural diversity but is semantically agnostic: the priors may not cover the natural image manifold in \textcolor{blue}{\uwave{complex}} or domain-specific tasks where intra-class variation is subtle. \textcolor{blue}{\uwave{In}} \textcolor{blue}{\uwave{particular,}} \textcolor{blue}{\uwave{semantic}} \textcolor{blue}{\uwave{cutmixing}} \textcolor{blue}{\uwave{uses}} \textcolor{blue}{\uwave{irregular}} \textcolor{blue}{\uwave{masks}} \textcolor{blue}{\uwave{rather}} \textcolor{blue}{\uwave{than}} \textcolor{blue}{\uwave{object}} \textcolor{blue}{\uwave{segments,}} \textcolor{blue}{\uwave{so}} \textcolor{blue}{\uwave{a}} \textcolor{blue}{\uwave{teacher}} \textcolor{blue}{\uwave{label}} \textcolor{blue}{\uwave{on}} \textcolor{blue}{\uwave{a}} \textcolor{blue}{\uwave{synthetic}} \textcolor{blue}{\uwave{image}} \textcolor{blue}{\uwave{does}} \textcolor{blue}{\uwave{not}} \textcolor{blue}{\uwave{make}} \textcolor{blue}{\uwave{it}} \textcolor{blue}{\uwave{a}} \textcolor{blue}{\uwave{meaningful}} \textcolor{blue}{\uwave{instance}} \textcolor{blue}{\uwave{of}} \textcolor{blue}{\uwave{that}} \textcolor{blue}{\uwave{class.}} & 23.75 & 121 \\
\addlinespace[3pt]
56 & \textit{(new text)} & \textcolor{blue}{\uwave{Each}} \textcolor{blue}{\uwave{method}} \textcolor{blue}{\uwave{pays}} \textcolor{blue}{\uwave{for}} \textcolor{blue}{\uwave{fidelity}} \textcolor{blue}{\uwave{with}} \textcolor{blue}{\uwave{extra}} \textcolor{blue}{\uwave{computation}} \textcolor{blue}{\uwave{or}} \textcolor{blue}{\uwave{teacher}} \textcolor{blue}{\uwave{access,}} \textcolor{blue}{\uwave{and}} \textcolor{blue}{\uwave{we}} \textcolor{blue}{\uwave{did}} \textcolor{blue}{\uwave{not}} \textcolor{blue}{\uwave{fully}} \textcolor{blue}{\uwave{measure}} \textcolor{blue}{\uwave{these}} \textcolor{blue}{\uwave{costs.}} \textcolor{blue}{\uwave{DivBFKD}} \textcolor{blue}{\uwave{queries}} \textcolor{blue}{\uwave{the}} \textcolor{blue}{\uwave{teacher}} \textcolor{blue}{\uwave{on}} \textcolor{blue}{\uwave{generated}} \textcolor{blue}{\uwave{images}} \textcolor{blue}{\uwave{throughout}} \textcolor{blue}{\uwave{adversarial}} \textcolor{blue}{\uwave{training}} \textcolor{blue}{\uwave{to}} \textcolor{blue}{\uwave{compute}} \textcolor{blue}{\uwave{confidence}} \textcolor{blue}{\uwave{scores,}} \textcolor{blue}{\uwave{and}} \textcolor{blue}{\uwave{again}} \textcolor{blue}{\uwave{to}} \textcolor{blue}{\uwave{label}} \textcolor{blue}{\uwave{the}} \textcolor{blue}{\uwave{synthetic}} \textcolor{blue}{\uwave{distillation}} \textcolor{blue}{\uwave{set.}} \textcolor{blue}{\uwave{DIPKD}} \textcolor{blue}{\uwave{needs}} \textcolor{blue}{\uwave{teacher}} \textcolor{blue}{\uwave{labels}} \textcolor{blue}{\uwave{for}} \textcolor{blue}{\uwave{its}} \textcolor{blue}{\uwave{synthetic}} \textcolor{blue}{\uwave{images}} \textcolor{blue}{\uwave{and}} \textcolor{blue}{\uwave{trains}} \textcolor{blue}{\uwave{a}} \textcolor{blue}{\uwave{primer}} \textcolor{blue}{\uwave{student}} \textcolor{blue}{\uwave{and}} \textcolor{blue}{\uwave{runs}} \textcolor{blue}{\uwave{a}} \textcolor{blue}{\uwave{contrastive}} \textcolor{blue}{\uwave{phase}} \textcolor{blue}{\uwave{before}} \textcolor{blue}{\uwave{the}} \textcolor{blue}{\uwave{final}} \textcolor{blue}{\uwave{student,}} \textcolor{blue}{\uwave{which}} \textcolor{blue}{\uwave{is}} \textcolor{blue}{\uwave{more}} \textcolor{blue}{\uwave{training}} \textcolor{blue}{\uwave{than}} \textcolor{blue}{\uwave{standard}} \textcolor{blue}{\uwave{distillation.}} \textcolor{blue}{\uwave{TAKE}} \textcolor{blue}{\uwave{adds}} \textcolor{blue}{\uwave{per-sample}} \textcolor{blue}{\uwave{gradient}} \textcolor{blue}{\uwave{tracking}} \textcolor{blue}{\uwave{over}} \textcolor{blue}{\uwave{the}} \textcolor{blue}{\uwave{training}} \textcolor{blue}{\uwave{trajectory,}} \textcolor{blue}{\uwave{LLM}} \textcolor{blue}{\uwave{generation,}} \textcolor{blue}{\uwave{and}} \textcolor{blue}{\uwave{an}} \textcolor{blue}{\uwave{optimal-transport}} \textcolor{blue}{\uwave{solve}} \textcolor{blue}{\uwave{whose}} \textcolor{blue}{\uwave{per-iteration}} \textcolor{blue}{\uwave{cost}} \textcolor{blue}{\uwave{grows}} \textcolor{blue}{\uwave{with}} \textcolor{blue}{\uwave{the}} \textcolor{blue}{\uwave{product}} \textcolor{blue}{\uwave{of}} \textcolor{blue}{\uwave{the}} \textcolor{blue}{\uwave{source}} \textcolor{blue}{\uwave{and}} \textcolor{blue}{\uwave{pool}} \textcolor{blue}{\uwave{sizes.}} \textcolor{blue}{\uwave{Whether}} \textcolor{blue}{\uwave{the}} \textcolor{blue}{\uwave{accuracy}} \textcolor{blue}{\uwave{gains}} \textcolor{blue}{\uwave{justify}} \textcolor{blue}{\uwave{these}} \textcolor{blue}{\uwave{costs}} \textcolor{blue}{\uwave{depends}} \textcolor{blue}{\uwave{on}} \textcolor{blue}{\uwave{how}} \textcolor{blue}{\uwave{expensive}} \textcolor{blue}{\uwave{teacher}} \textcolor{blue}{\uwave{queries}} \textcolor{blue}{\uwave{and}} \textcolor{blue}{\uwave{real}} \textcolor{blue}{\uwave{data}} \textcolor{blue}{\uwave{are}} \textcolor{blue}{\uwave{in}} \textcolor{blue}{\uwave{a}} \textcolor{blue}{\uwave{given}} \textcolor{blue}{\uwave{deployment.}} & 87.50 & 122 \\
\addlinespace[3pt]
57 & \textit{(new text)} & \textcolor{blue}{\uwave{The}} \textcolor{blue}{\uwave{evidence}} \textcolor{blue}{\uwave{for}} \textcolor{blue}{\uwave{generalization}} \textcolor{blue}{\uwave{is}} \textcolor{blue}{\uwave{bounded}} \textcolor{blue}{\uwave{by}} \textcolor{blue}{\uwave{what}} \textcolor{blue}{\uwave{we}} \textcolor{blue}{\uwave{tested.}} \textcolor{blue}{\uwave{Claims}} \textcolor{blue}{\uwave{of}} \textcolor{blue}{\uwave{generality}} \textcolor{blue}{\uwave{across}} \textcolor{blue}{\uwave{architectures,}} \textcolor{blue}{\uwave{datasets,}} \textcolor{blue}{\uwave{or}} \textcolor{blue}{\uwave{domains}} \textcolor{blue}{\uwave{should}} \textcolor{blue}{\uwave{be}} \textcolor{blue}{\uwave{read}} \textcolor{blue}{\uwave{as}} \textcolor{blue}{\uwave{limited}} \textcolor{blue}{\uwave{to}} \textcolor{blue}{\uwave{the}} \textcolor{blue}{\uwave{settings}} \textcolor{blue}{\uwave{evaluated}} \textcolor{blue}{\uwave{in}} \textcolor{blue}{\uwave{each}} \textcolor{blue}{\uwave{chapter,}} \textcolor{blue}{\uwave{and}} \textcolor{blue}{\uwave{the}} \textcolor{blue}{\uwave{theoretical}} \textcolor{blue}{\uwave{guarantees}} \textcolor{blue}{\uwave{hold}} \textcolor{blue}{\uwave{only}} \textcolor{blue}{\uwave{under}} \textcolor{blue}{\uwave{their}} \textcolor{blue}{\uwave{stated}} \textcolor{blue}{\uwave{assumptions.}} & 25.75 & 122 \\
\addlinespace[3pt]
\multicolumn{5}{l}{\textbf{6.2 Limitations and Future Directions $\rightarrow$ 6.2.2 Theoretical Open Problems}}\\[2pt]
58 & No formal bound relates the distributional coverage of a synthetic or selected set to the generalization gap of a student trained on it, because the relevant bound must account jointly for the synthesis or selection procedure, the student architecture, and the task loss. & \textcolor{blue}{\uwave{Beyond}} \textcolor{blue}{\uwave{the}} \textcolor{blue}{\uwave{assumption-dependent}} \textcolor{blue}{\uwave{loss-gap}} \textcolor{blue}{\uwave{bound}} \textcolor{blue}{\uwave{for}} \textcolor{blue}{\uwave{TAKE}} \textcolor{blue}{\uwave{(Theorem}} \textcolor{blue}{\uwave{\#),}} no \textcolor{blue}{\uwave{general}} formal bound relates the distributional coverage of a synthetic or selected set to the generalization gap of a student trained on it, because the relevant bound must account jointly for the synthesis or selection procedure, the student architecture, and the task loss. \textcolor{blue}{\uwave{Relatedly,}} \textcolor{blue}{\uwave{distributional}} \textcolor{blue}{\uwave{fidelity}} \textcolor{blue}{\uwave{is}} \textcolor{blue}{\uwave{a}} \textcolor{blue}{\uwave{design}} \textcolor{blue}{\uwave{principle}} \textcolor{blue}{\uwave{that}} \textcolor{blue}{\uwave{the}} \textcolor{blue}{\uwave{three}} \textcolor{blue}{\uwave{methods}} \textcolor{blue}{\uwave{operationalize}} \textcolor{blue}{\uwave{differently,}} \textcolor{blue}{\uwave{not}} \textcolor{blue}{\uwave{a}} \textcolor{blue}{\uwave{general}} \textcolor{blue}{\uwave{theoretical}} \textcolor{blue}{\uwave{framework,}} \textcolor{blue}{\uwave{and}} \textcolor{blue}{\uwave{no}} \textcolor{blue}{\uwave{common}} \textcolor{blue}{\uwave{quantitative}} \textcolor{blue}{\uwave{measure}} \textcolor{blue}{\uwave{yet}} \textcolor{blue}{\uwave{links}} \textcolor{blue}{\uwave{it}} \textcolor{blue}{\uwave{to}} \textcolor{blue}{\uwave{downstream}} \textcolor{blue}{\uwave{performance}} \textcolor{blue}{\uwave{across}} \textcolor{blue}{\uwave{the}} \textcolor{blue}{\uwave{three}} \textcolor{blue}{\uwave{settings.}} & 34.00 & 122 \\
\addlinespace[3pt]
\end{longtable}
\end{document}
